% ============================================================================
% Préambule --- Rapport de mini-projet « Jumeau numérique d'une ligne de production »
% ENSAM Rabat --- Filière INDIA&SD --- Industrie 4.0 --- Année Universitaire 2026/2027
% Gabarit repris du rapport de stage PIPER Speech (même charte : bleu marine ENSAM).
% ============================================================================

\usepackage[utf8]{inputenc}
\usepackage[T1]{fontenc}
\usepackage[french]{babel}
\usepackage{lmodern}
\usepackage{geometry}
\usepackage{graphicx}
\usepackage{listings}
\usepackage{xcolor}
\usepackage{hyperref}
\usepackage{fancyhdr}
\usepackage{titlesec}
\usepackage{booktabs}
\usepackage{tabularx}
\usepackage{longtable}
\usepackage{enumitem}
\usepackage{tikz}
\usepackage{float}
\usepackage{amsmath}
\usepackage{amssymb}
\usepackage{caption}
\usepackage{subcaption}
\usepackage{tcolorbox}
\usepackage{acronym}
\usepackage{csquotes}
\usepackage{colortbl}
\usepackage{multirow}
\usepackage{pgfplots}
\pgfplotsset{compat=1.18}

\usetikzlibrary{shapes.geometric, arrows.meta, positioning, fit, backgrounds, calc}

% ---------------------------- Mise en page ----------------------------------
\geometry{a4paper, margin=2.5cm, headheight=15pt}

% Couleur principale : bleu marine académique ENSAM (identique au rapport PIPER),
% couleur secondaire : rouge du logo ENSAM.
\definecolor{ensambleu}{HTML}{1B3A5C}
\definecolor{ensamrouge}{HTML}{B22A2E}
\definecolor{lienviolet}{HTML}{4A148C}
\definecolor{etatmarche}{HTML}{22C55E}
\definecolor{etatattente}{HTML}{94A3B8}
\definecolor{etatbloque}{HTML}{F59E0B}
\definecolor{etatpanne}{HTML}{EF4444}

\hypersetup{
    colorlinks=true,
    linkcolor=lienviolet,
    urlcolor=lienviolet,
    citecolor=lienviolet,
    pdftitle={Jumeau numérique d'une ligne de production},
    pdfauthor={Issami, Ezbakhe, Benzahra, Hsain, Kichchou, Hamdouch}
}

% ---------------------------- Listings (code) -------------------------------
\definecolor{codebg}{HTML}{F5F5F5}
\definecolor{codegreen}{HTML}{28A745}
\definecolor{codepurple}{HTML}{6F42C1}
\definecolor{codegray}{HTML}{6A737D}

\lstdefinestyle{mystyle}{
    backgroundcolor=\color{codebg},
    commentstyle=\color{codegray}\itshape,
    keywordstyle=\color{codepurple}\bfseries,
    numberstyle=\tiny\color{codegray},
    stringstyle=\color{codegreen},
    basicstyle=\ttfamily\footnotesize,
    breakatwhitespace=false,
    breaklines=true,
    captionpos=b,
    keepspaces=true,
    numbers=left,
    numbersep=5pt,
    showspaces=false,
    showstringspaces=false,
    showtabs=false,
    tabsize=2,
    frame=single,
    rulecolor=\color{gray!30},
    extendedchars=true,
    inputencoding=utf8,
    literate=%
      {é}{{\'e}}1 {è}{{\`e}}1 {ê}{{\^e}}1 {ë}{{\"e}}1
      {à}{{\`a}}1 {â}{{\^a}}1 {ä}{{\"a}}1
      {î}{{\^i}}1 {ï}{{\"i}}1
      {ô}{{\^o}}1 {ö}{{\"o}}1
      {ù}{{\`u}}1 {û}{{\^u}}1 {ü}{{\"u}}1
      {ç}{{\c c}}1 {Ç}{{\c C}}1
      {É}{{\'E}}1 {È}{{\`E}}1 {À}{{\`A}}1 {×}{{$\times$}}1
}
\lstset{style=mystyle}

% ---------------------------- En-tête / pied --------------------------------
\renewcommand{\chaptermark}[1]{\markboth{\chaptername\ \thechapter}{}}

\pagestyle{fancy}
\fancyhf{}
\fancyhead[L]{\small\nouppercase{\leftmark}}
\fancyhead[R]{\small Jumeau numérique --- Industrie 4.0}
\fancyfoot[L]{\footnotesize ENSAM Rabat --- Année Universitaire 2026/2027}
\fancyfoot[R]{\thepage}
\renewcommand{\headrulewidth}{0.4pt}
\renewcommand{\footrulewidth}{0.4pt}

\fancypagestyle{plain}{%
  \fancyhf{}%
  \fancyfoot[L]{\footnotesize ENSAM Rabat --- Année Universitaire 2026/2027}%
  \fancyfoot[R]{\thepage}%
  \renewcommand{\headrulewidth}{0pt}%
  \renewcommand{\footrulewidth}{0.4pt}%
}

% ---------------------------- Titres chapitres ------------------------------
\titleformat{\chapter}[display]
  {\normalfont\huge\bfseries\color{ensambleu}}
  {\chaptertitlename\ \thechapter}{20pt}{\Huge}
\titlespacing*{\chapter}{0pt}{-20pt}{30pt}

% ---------------------------- Légendes --------------------------------------
\captionsetup{font=small, labelfont=bf, labelsep=period}

% ---------------------------- Macros utiles ---------------------------------
\newcommand{\produit}{\textbf{jumeau numérique}}
\newcommand{\code}[1]{\texttt{#1}}

\newtcolorbox{callout}[1][]{colback=ensambleu!5, colframe=ensambleu, sharp corners,
  boxrule=0.5pt, left=8pt, right=8pt, top=6pt, bottom=6pt, #1}
\newtcolorbox{success}[1][]{colback=green!6, colframe=green!60!black, sharp corners,
  boxrule=0.5pt, left=8pt, right=8pt, top=6pt, bottom=6pt, #1}
\newtcolorbox{warn}[1][]{colback=orange!6, colframe=orange!70!black, sharp corners,
  boxrule=0.5pt, left=8pt, right=8pt, top=6pt, bottom=6pt, #1}

\setlist[itemize]{itemsep=2pt, topsep=3pt}
\setlist[enumerate]{itemsep=2pt, topsep=3pt}

% ---------------------------- Acronymes -------------------------------------
\acrodef{3D}{Trois dimensions}
\acrodef{API}{Application Programming Interface}
\acrodef{CSV}{Comma-Separated Values}
\acrodef{ENSAM}{École Nationale Supérieure d'Arts et Métiers}
\acrodef{ERP}{Enterprise Resource Planning}
\acrodef{IHM}{Interface Homme-Machine}
\acrodef{INDIA}[INDIA\&SD]{Ingénierie Numérique en Data Science, Intelligence Artificielle \& Santé Digitale}
\acrodef{IoT}{Internet of Things}
\acrodef{JSON}{JavaScript Object Notation}
\acrodef{KPI}{Key Performance Indicator}
\acrodef{MES}{Manufacturing Execution System}
\acrodef{MTBF}{Mean Time Between Failures}
\acrodef{MTTR}{Mean Time To Repair}
\acrodef{OPC}[OPC UA]{Open Platform Communications Unified Architecture}
\acrodef{PLC}{Programmable Logic Controller}
\acrodef{REST}{Representational State Transfer}
\acrodef{SQL}{Structured Query Language}
\acrodef{TRS}{Taux de Rendement Synthétique}
\acrodef{WIP}{Work In Progress (en-cours)}
\acrodef{WS}{WebSocket}

\endinput
